\subsection{مبانی نظری و فرمول‌بندی ریاضی}
فرآیند فشرده‌سازی در \lr{Bonsai Pruner} بر پایه اصول ریاضی زیر فرمول‌بندی شده است:
\begin{enumerate}
    \item \textbf{زمان‌بندی مکعبی تنکی (\lr{Cubic Schedule}):} برای دستیابی به نسبت تنکی نهایی $S_{\text{target}}$ در طول $T$ گام متوالی، تنکی هدف در گام $t \in \{1, \dots, T\}$ طبق رابطه زیر تعیین می‌شود:
    \begin{equation}
    S_t = S_{\text{target}} \left( 1 - \left( 1 - \frac{t}{T} \right)^3 \right)
    \end{equation}
    این رابطه تضمین می‌کند که در گام‌های آغازین نرخ حذف وزن‌ها بسیار آرام و کنترل‌شده بوده و در گام‌های پایانی به آرامی به سطح هدف همگرا گردد.
    \item \textbf{معیار حساسیت تیلور مرتبه اول با حفاظت \lr{[CLS]}:} در هر گام، اهمیت هر پارامتر وزن بر اساس ضرب اندازه وزن در گرادیان تجربی سنجیده می‌شود:
    \begin{equation}
    \mathbf{S}_{ij} = |\mathbf{W}_{ij}| \cdot \left| \frac{\partial \mathcal{L}}{\partial \mathbf{W}_{ij}} \right|
    \end{equation}
    در لایه‌های توجه، سطرها و ستون‌های مرتبط با هدایت توکن \lr{[CLS]} با ضریب حفاظتی $\kappa > 1$ تقویت می‌گردند تا در اولویت حذفیات قرار نگیرند: $\mathbf{S}_{\text{CLS}} \leftarrow \kappa \cdot \mathbf{S}_{\text{CLS}}$.
    \item \textbf{گام‌های بازیابی گرادیانی بین‌مرحله‌ای (\lr{Gradient Recovery}):} پس از به‌روزرسانی ماسک تنک $\mathbf{M}_t$، وزن‌های باقیمانده به تعداد $K$ تکرار ریزگرادیانی با نرخ یادگیری کوچک $\eta_{\text{rec}}$ طبق رابطه زیر بهینه‌سازی می‌شوند تا انحراف لاجیت‌ها جبران گردد:
    \begin{equation}
    \mathbf{W} \leftarrow \mathbf{W} - \eta_{\text{rec}} \cdot \left( \nabla_{\mathbf{W}} \mathcal{L}_{\text{task}} \odot \mathbf{M}_t \right)
    \end{equation}
\end{enumerate}
